\documentclass[cs4size,a4paper]{ctexart}   
%==================== 数学符号公式 ============
\usepackage{caption}
\usepackage{subfigure}
\usepackage[algo2e,ruled,vlined,linesnumbered]{algorithm2e} % 伪代码包
\usepackage{enumerate}
\usepackage{amsmath}                 % AMS LaTeX宏包
\usepackage[style=1]{mdframed}
\usepackage{amsthm}
\usepackage{amsfonts}
\usepackage{mathrsfs}                % 英文花体字 体
\usepackage{bm}                      % 数学公式中的黑斜体
\usepackage{bbding,manfnt}           % 一些图标，如 \dbend
\usepackage{lettrine}                % 首字下沉，命令\lettrine
\def\attention{\lettrine[lines=2,lraise=0,nindent=0em]{\large\textdbend\hspace{1mm}}{}}
\usepackage{longtable}
\usepackage[toc,page]{appendix}
\usepackage{geometry}                % 页边距调整
\geometry{top=3.0cm,bottom=2.7cm,left=2.5cm,right=2.5cm}
%====================公式按章编号==========================
\numberwithin{equation}{section}
\numberwithin{table}{section}
\numberwithin{figure}{section}
%================= 基本格式预置 ===========================
\usepackage{fancyhdr}
\pagestyle{fancy}
\fancyhf{}  
\fancyhead[C]{\zihao{5}  \kaishu 机\ 器\ 学\ 习\ 结\ 课\ 报\ 告}
\fancyfoot[C]{~\zihao{5} \thepage~}
\renewcommand{\headrulewidth}{0.65pt} 
%\CTEXsetup[format={\centering\bfseries\zihao{-2}}]{section}
\ctexset{
    section={format={\centering\bfseries\zihao{-2}}},
    subsection={nameformat={\bfseries\zihao{3}}},
    subsubsection={nameformat={\bfseries\zihao{4}}}
}
%================== 图形支持宏包 =========================
\usepackage{subfigure}
\usepackage{graphicx}                % 嵌入png图像
\usepackage{color,xcolor}            % 支持彩色文本、底色、文本框等
\usepackage{hyperref}                % 交叉引用
\usepackage{caption}
\captionsetup{figurewithin=section}
%==================== 源码和流程图 =====================
\usepackage{listings}                % 粘贴源代码
\usepackage{xcolor}
\usepackage{color}
\definecolor{dkgreen}{rgb}{0,0.6,0}
\definecolor{gray}{rgb}{0.5,0.5,0.5}
\definecolor{mauve}{rgb}{0.58,0,0.82}
\usepackage{xcolor}
\lstset{
    %行号
    numbers=left,
    %背景框
    framexleftmargin=8mm,
    frame=none,
    %背景色
    %backgroundcolor=\color[rgb]{1,1,0.76},
    backgroundcolor=\color[RGB]{245,245,244},
    %样式
    keywordstyle=\bf\color{blue},
    identifierstyle=\bf,
    numberstyle=\color[RGB]{0,192,192},
    commentstyle=\it\color[RGB]{0,96,96},
    stringstyle=\rmfamily\slshape\color[RGB]{128,0,0},
    %显示空格
    showstringspaces=false
}


%--------------------
\hypersetup{hidelinks}
\usepackage{booktabs}  
\usepackage{shorttoc}
\usepackage{tabu,tikz}
\usepackage{float}

\usepackage{multirow}



\tabcolsep=1ex
\tabulinesep=\tabcolsep
\newlength\tikzboxwidth
\newlength\tikzboxheight
\newcommand\tikzbox[1]{
    \settowidth\tikzboxwidth{#1}
    \settoheight\tikzboxheight{#1}
    \begin{tikzpicture}
        \path[use as bounding box]
            (-0.5\tikzboxwidth,-0.5\tikzboxheight)rectangle
            (0.5\tikzboxwidth,0.5\tikzboxheight);
        \node[inner sep=\tabcolsep+0.5\arrayrulewidth,line width=0.5mm,draw=black]
            at(0,0){#1};
    \end{tikzpicture}%
}

\makeatletter
\def\hlinew#1{
    \noalign{\ifnum0=`}\fi\hrule \@height #1 \futurelet
    \reserved@a\@xhline
}
   
\newcommand{\tabincell}[2]{\begin{tabular}{@{}#1@{}}#2\end{tabular}}%

\usepackage{subfigure}

\usepackage{ifthen}


\usepackage{graphicx} 
\newcommand{\HRule}{\rule{\linewidth}{0.5mm}}

\newtheorem{Theorem}{定理}
\newtheorem{Lemma}{引理} 
%%使得公式随章节自动编号
\makeatletter
\@addtoreset{equation}{section}
\makeatother
\renewcommand{\theequation}{\arabic{section}.\arabic{equation}}

%-------------------------
	
\usepackage{pythonhighlight}
\usepackage{tikz}                    
\usepackage{tikz-3dplot}
\usetikzlibrary{shapes,arrows,positioning}
%===================   正文开始    ===================
\begin{document}
\bibliographystyle{gbt7714-2005}     %论文引用格式
%===================  定理类环境定义 ===================
\newtheorem{example}{例}              % 整体编号
\newtheorem{algorithm}{算法}
\newtheorem{theorem}{定理}            % 按 section 编号
\newtheorem{definition}{定义}
\newtheorem{axiom}{公理}
\newtheorem{property}{性质}
\newtheorem{proposition}{命题}
\newtheorem{lemma}{引理}
\newtheorem{corollary}{推论}
\newtheorem{remark}{注解}
\newtheorem{condition}{条件}
\newtheorem{conclusion}{结论}
\newtheorem{assumption}{假设}
%==================重定义 ===================
\renewcommand{\contentsname}{目录}     
\renewcommand{\abstractname}{摘要} 
\renewcommand{\refname}{参考文献}     
\renewcommand{\indexname}{索引}
\renewcommand{\figurename}{图}
\renewcommand{\tablename}{表}
\renewcommand{\appendixname}{附录}
\renewcommand{\proofname}{证明}
\renewcommand{\algorithm}{算法} 
%============== 封皮和前言 =================
\begin{titlepage}

\begin{center}


% Upper part of the page
 

{\huge \bfseries 华中科技大学计算机科学与技术学院}\\[1.5cm]

{ \huge \bfseries 《机器学习》结课报告}\\[1.5cm]

\includegraphics[width=0.55\textwidth]{figure/logo1.png}\\[1.5cm]   



% Title


% Author and supervisor

\begin{minipage}{0.5\textwidth}
\large
{\heiti 专~~~~~~~业}\quad \underline{~~~~~~~~~~~~~~~~~~~~~~~~~~~}

{\heiti 班~~~~~~~级}\quad \underline{~~~~~~~~~~~~~~~~~~~~~~~~~~~}

{\heiti 学~~~~~~~号}\quad \underline{~~~~~~~~~~~~~~~~~~~~~~~~~~~}

{\heiti 姓~~~~~~~名}\quad \underline{~~~~~~~~~~~~~~~~~~~~~~~~~~~}

{\heiti 成~~~~~~~绩}\quad \underline{~~~~~~~~~~~~~~~~~~~~~~~~~~~}

{\heiti 指导教师}\quad \underline{~~~~~~~~~~~~~~~~~~~~~~~~~~~}

{\heiti 时~~~~~~~间}\quad \underline{\today~~}
\end{minipage}

\vfill

\end{center}

\end{titlepage}

\pagestyle{plain}
\pagenumbering{Roman}
\pagestyle{empty}
\tableofcontents 
\thispagestyle{empty}
%============== 论文正文   =================
\pagestyle{fancy}
\pagenumbering{arabic}
\setlength{\headheight}{12.64723pt}
\addtolength{\topmargin}{-0.64723pt}
\section{实验要求}

总体要求：

\begin{enumerate}
    \item 控制报告页数，不要大段大段贴代码
    \item 表格和插图请编号并进行交叉引用，表格使用三线表
\end{enumerate}







\section{算法设计与实现}



\section{实验环境与平台}






\section{结果与分析}




\section{个人体会}



%============= 参考文献 =====================
\addcontentsline{toc}{section}{参考文献}
\begin{thebibliography}{99}  


\bibitem{ref1} 周志华. 机器学习:第3章. 清华大学出版社, 2016.



\end{thebibliography}
\clearpage
%=============  致谢  ======================

\end{document}
%%%%%%%%%% 结束 %%%%%%%%%%